% ============================================================
%  main.tex  —  Demo presentation using the HPC Beamer theme
%  Aspect ratio 16:9.  Compile with: pdflatex main.tex (×2)
% ============================================================

\documentclass[aspectratio=169]{beamer}

% ---- Packages -----------------------------------------------
\usepackage[utf8]{inputenc}
\usepackage[T1]{fontenc}
\usepackage{amsmath, amssymb, amsthm}
\usepackage{mathtools}
\usepackage{graphicx}
\usepackage{booktabs}       % professional tables
\usepackage{multicol}
\usepackage{listings}       % code listings
\usepackage{hyperref}
\usepackage[yyyymmdd]{datetime}
\renewcommand{\dateseparator}{-}

% ---- Load theme (split architecture: 5 .sty files) ----------
\usetheme[watermark=on, logo=none]{horixon}

% ---- Load environment shortcuts + user macros ---------------
% ============================================================
%  environment-shortcuts.tex
%  Convenience wrappers for common Beamer environments.
%  % ============================================================
%  environment-shortcuts.tex
%  Convenience wrappers for common Beamer environments.
%  % ============================================================
%  environment-shortcuts.tex
%  Convenience wrappers for common Beamer environments.
%  \input{environment-shortcuts} in your preamble.
%
%  ADDING YOUR OWN MACROS:
%  Place them after the "── User macros ──" divider below.
% ============================================================


% ---- Layout helpers -----------------------------------------
\newcommand{\columnsbegin}{\begin{columns}}
\newcommand{\columnsend}{\end{columns}}
\newcommand{\colbegin}[1]{\begin{column}{#1}}
\newcommand{\colend}{\end{column}}


% ---- General content blocks ---------------------------------
\newcommand{\blockbegin}[1]{\begin{block}{#1}}
\newcommand{\blockend}{\end{block}}

\newcommand{\alertblockbegin}[1]{\begin{alertblock}{#1}}
\newcommand{\alertblockend}{\end{alertblock}}

\newcommand{\exampleblockbegin}[1]{\begin{exampleblock}{#1}}
\newcommand{\exampleblockend}{\end{exampleblock}}


% ---- Mathematical environments ------------------------------
\newcommand{\definitionbegin}[1]{\begin{definition}[#1]}
\newcommand{\definitionend}{\end{definition}}

\newcommand{\theorembegin}[1]{\begin{theorem}[#1]}
\newcommand{\theoremend}{\end{theorem}}

\newcommand{\lemmabegin}[1]{\begin{lemma}[#1]}
\newcommand{\lemmaend}{\end{lemma}}

\newcommand{\corollarybegin}[1]{\begin{corollary}[#1]}
\newcommand{\corollaryend}{\end{corollary}}

\newcommand{\proofbegin}{\begin{proof}}
\newcommand{\proofend}{\end{proof}}

\newcommand{\examplebegin}[1]{\begin{example}[#1]}
\newcommand{\exampleend}{\end{example}}

\newcommand{\remarkbegin}[1]{\begin{remark}[#1]}
\newcommand{\remarkend}{\end{remark}}


% ---- Number-set shortcuts (common in scientific writing) ----
\newcommand{\R}{\mathbb{R}}   % Reals
\newcommand{\Z}{\mathbb{Z}}   % Integers
\newcommand{\N}{\mathbb{N}}   % Naturals
\newcommand{\C}{\mathbb{C}}   % Complex
\newcommand{\Q}{\mathbb{Q}}   % Rationals


% ---- HPC / numerics shorthands ------------------------------
\newcommand{\BigO}[1]{\ensuremath{\mathcal{O}\!\left(#1\right)}}
\newcommand{\norm}[1]{\ensuremath{\left\lVert#1\right\rVert}}
\newcommand{\abs}[1]{\ensuremath{\left|#1\right|}}
\newcommand{\inner}[2]{\ensuremath{\langle #1,\, #2 \rangle}}
\newcommand{\vect}[1]{\ensuremath{\boldsymbol{#1}}}
\newcommand{\mat}[1]{\ensuremath{\mathbf{#1}}}
\newcommand{\T}{\ensuremath{^{\top}}}


% ================================================================
%  ── User macros ─────────────────────────────────────────────
%  Add your own \newcommand / \renewcommand definitions below.
% ================================================================


 in your preamble.
%
%  ADDING YOUR OWN MACROS:
%  Place them after the "── User macros ──" divider below.
% ============================================================


% ---- Layout helpers -----------------------------------------
\newcommand{\columnsbegin}{\begin{columns}}
\newcommand{\columnsend}{\end{columns}}
\newcommand{\colbegin}[1]{\begin{column}{#1}}
\newcommand{\colend}{\end{column}}


% ---- General content blocks ---------------------------------
\newcommand{\blockbegin}[1]{\begin{block}{#1}}
\newcommand{\blockend}{\end{block}}

\newcommand{\alertblockbegin}[1]{\begin{alertblock}{#1}}
\newcommand{\alertblockend}{\end{alertblock}}

\newcommand{\exampleblockbegin}[1]{\begin{exampleblock}{#1}}
\newcommand{\exampleblockend}{\end{exampleblock}}


% ---- Mathematical environments ------------------------------
\newcommand{\definitionbegin}[1]{\begin{definition}[#1]}
\newcommand{\definitionend}{\end{definition}}

\newcommand{\theorembegin}[1]{\begin{theorem}[#1]}
\newcommand{\theoremend}{\end{theorem}}

\newcommand{\lemmabegin}[1]{\begin{lemma}[#1]}
\newcommand{\lemmaend}{\end{lemma}}

\newcommand{\corollarybegin}[1]{\begin{corollary}[#1]}
\newcommand{\corollaryend}{\end{corollary}}

\newcommand{\proofbegin}{\begin{proof}}
\newcommand{\proofend}{\end{proof}}

\newcommand{\examplebegin}[1]{\begin{example}[#1]}
\newcommand{\exampleend}{\end{example}}

\newcommand{\remarkbegin}[1]{\begin{remark}[#1]}
\newcommand{\remarkend}{\end{remark}}


% ---- Number-set shortcuts (common in scientific writing) ----
\newcommand{\R}{\mathbb{R}}   % Reals
\newcommand{\Z}{\mathbb{Z}}   % Integers
\newcommand{\N}{\mathbb{N}}   % Naturals
\newcommand{\C}{\mathbb{C}}   % Complex
\newcommand{\Q}{\mathbb{Q}}   % Rationals


% ---- HPC / numerics shorthands ------------------------------
\newcommand{\BigO}[1]{\ensuremath{\mathcal{O}\!\left(#1\right)}}
\newcommand{\norm}[1]{\ensuremath{\left\lVert#1\right\rVert}}
\newcommand{\abs}[1]{\ensuremath{\left|#1\right|}}
\newcommand{\inner}[2]{\ensuremath{\langle #1,\, #2 \rangle}}
\newcommand{\vect}[1]{\ensuremath{\boldsymbol{#1}}}
\newcommand{\mat}[1]{\ensuremath{\mathbf{#1}}}
\newcommand{\T}{\ensuremath{^{\top}}}


% ================================================================
%  ── User macros ─────────────────────────────────────────────
%  Add your own \newcommand / \renewcommand definitions below.
% ================================================================


 in your preamble.
%
%  ADDING YOUR OWN MACROS:
%  Place them after the "── User macros ──" divider below.
% ============================================================


% ---- Layout helpers -----------------------------------------
\newcommand{\columnsbegin}{\begin{columns}}
\newcommand{\columnsend}{\end{columns}}
\newcommand{\colbegin}[1]{\begin{column}{#1}}
\newcommand{\colend}{\end{column}}


% ---- General content blocks ---------------------------------
\newcommand{\blockbegin}[1]{\begin{block}{#1}}
\newcommand{\blockend}{\end{block}}

\newcommand{\alertblockbegin}[1]{\begin{alertblock}{#1}}
\newcommand{\alertblockend}{\end{alertblock}}

\newcommand{\exampleblockbegin}[1]{\begin{exampleblock}{#1}}
\newcommand{\exampleblockend}{\end{exampleblock}}


% ---- Mathematical environments ------------------------------
\newcommand{\definitionbegin}[1]{\begin{definition}[#1]}
\newcommand{\definitionend}{\end{definition}}

\newcommand{\theorembegin}[1]{\begin{theorem}[#1]}
\newcommand{\theoremend}{\end{theorem}}

\newcommand{\lemmabegin}[1]{\begin{lemma}[#1]}
\newcommand{\lemmaend}{\end{lemma}}

\newcommand{\corollarybegin}[1]{\begin{corollary}[#1]}
\newcommand{\corollaryend}{\end{corollary}}

\newcommand{\proofbegin}{\begin{proof}}
\newcommand{\proofend}{\end{proof}}

\newcommand{\examplebegin}[1]{\begin{example}[#1]}
\newcommand{\exampleend}{\end{example}}

\newcommand{\remarkbegin}[1]{\begin{remark}[#1]}
\newcommand{\remarkend}{\end{remark}}


% ---- Number-set shortcuts (common in scientific writing) ----
\newcommand{\R}{\mathbb{R}}   % Reals
\newcommand{\Z}{\mathbb{Z}}   % Integers
\newcommand{\N}{\mathbb{N}}   % Naturals
\newcommand{\C}{\mathbb{C}}   % Complex
\newcommand{\Q}{\mathbb{Q}}   % Rationals


% ---- HPC / numerics shorthands ------------------------------
\newcommand{\BigO}[1]{\ensuremath{\mathcal{O}\!\left(#1\right)}}
\newcommand{\norm}[1]{\ensuremath{\left\lVert#1\right\rVert}}
\newcommand{\abs}[1]{\ensuremath{\left|#1\right|}}
\newcommand{\inner}[2]{\ensuremath{\langle #1,\, #2 \rangle}}
\newcommand{\vect}[1]{\ensuremath{\boldsymbol{#1}}}
\newcommand{\mat}[1]{\ensuremath{\mathbf{#1}}}
\newcommand{\T}{\ensuremath{^{\top}}}


% ================================================================
%  ── User macros ─────────────────────────────────────────────
%  Add your own \newcommand / \renewcommand definitions below.
% ================================================================




% ---- Listings style (for code slides) -----------------------
\lstset{
  basicstyle=\ttfamily\scriptsize,
  keywordstyle=\color{horixonaccent}\bfseries,
  commentstyle=\color{horixonaccent!60!horixonnavy}\itshape,
  stringstyle=\color{horixondark},
  backgroundcolor=\color{horixondark!25!horixoncream},
  frame=single,
  framerule=0.4pt,
  rulecolor=\color{horixonaccent!70!white},
  numbers=left,
  numberstyle=\tiny\color{horixonaccent},
  breaklines=true,
  tabsize=2,
}

% ---- Meta information ---------------------------------------
\title[HPC Beamer Theme]{High-Performance Computing\\[0.1em]Beamer Presentation Theme}
\subtitle{A Custom Split-Architecture Template}
\author[Your Name]{Your Name}
\institute[Your Institution]{Department of Mathematics / Computer Science\\Your University}
\date{\today}

% =============================================================
\begin{document}
% =============================================================

% ---- Title page ---------------------------------------------
\begin{frame}[plain]
  \titlepage
\end{frame}

% ---- Table of contents --------------------------------------
\begin{frame}[plain]{Table of Contents}
  \tableofcontents
\end{frame}

% =============================================================
\section{Theme Architecture}
% =============================================================

\begin{frame}{Split Theme Architecture}
  This theme follows Beamer's recommended \emph{split-theme} pattern.
  Each concern lives in its own \texttt{.sty} file:

  \vspace{0.3em}
  \begin{columns}[t]
    \begin{column}{0.48\textwidth}
      \blockbegin{Style Files}
        \begin{itemize}
          \item \texttt{beamerthemehorixon.sty} — master loader
          \item \texttt{beamercolorthemehorixon.sty} — palette
          \item \texttt{beamerfontthemehorixon.sty} — typography
        \end{itemize}
      \blockend
    \end{column}
    \begin{column}{0.48\textwidth}
      \blockbegin{Layout Files}
        \begin{itemize}
          \item \texttt{beamerinnerthemehorixon.sty} — title page, blocks, bullets
          \item \texttt{beamerouterthemehorixon.sty} — header, footer, background
          \item \texttt{environment-shortcuts.tex} — user macros
        \end{itemize}
      \blockend
    \end{column}
  \end{columns}
\end{frame}

% =============================================================
\section{Color \& Font Palette}
% =============================================================

\begin{frame}{Color Palette}
  \begin{columns}[c]
    \begin{column}{0.55\textwidth}
      All colors are defined in \texttt{beamercolorthemehorixon.sty}
      and can be freely referenced in your slides:

      \vspace{0.4em}
      \begin{itemize}
        \item \textcolor{horixonnavy}{\rule{0.8em}{0.8em}}~\texttt{horixonnavy}   \texttt{\#222831} — primary background
        \item \textcolor{horixondark}{\rule{0.8em}{0.8em}}~\texttt{horixondark}   \texttt{\#393E46} — panels, header, footer
        \item \textcolor{horixonaccent}{\rule{0.8em}{0.8em}}~\texttt{horixonaccent} \texttt{\#948979} — accent, rules, bullets
        \item \textcolor{horixoncream}{\rule{0.8em}{0.8em}}~\texttt{horixoncream}  \texttt{\#DFD0B8} — text, light fills
        \item \textcolor{horixonalert}{\rule{0.8em}{0.8em}}~\texttt{horixonalert}  — alerts, warnings
      \end{itemize}
    \end{column}
    \begin{column}{0.40\textwidth}
      \blockbegin{Standard Block}
        Uses \texttt{horixoncyan} header.
      \blockend
      \alertblockbegin{Alert Block}
        Uses \texttt{horixonalert} header.
      \alertblockend
      \exampleblockbegin{Example Block}
        Uses \texttt{horixonlightnavy} header.
      \exampleblockend
    \end{column}
  \end{columns}
\end{frame}

% =============================================================
\section{Math \& HPC Macros}
% =============================================================

\begin{frame}{Mathematical Shortcuts}
  The shortcuts file provides common math macros out of the box.

  \vspace{0.4em}
  \theorembegin{Convergence of Conjugate Gradient}
    For $\mat{A} \in \R^{n \times n}$ symmetric positive definite,
    CG satisfies
    \[
      \norm{\vect{e}_k}_{\mat{A}}
        \;\leq\;
      2\left(\frac{\sqrt{\kappa}-1}{\sqrt{\kappa}+1}\right)^{\!k}
      \norm{\vect{e}_0}_{\mat{A}},
    \]
    where $\kappa = \lambda_{\max}/\lambda_{\min}$ is the condition number.
  \theoremend

  \vspace{0.3em}
  \begin{itemize}
    \item \texttt{\textbackslash vect\{x\}} $\to \vect{x}$ \quad
          \texttt{\textbackslash mat\{A\}}  $\to \mat{A}$ \quad
          \texttt{\textbackslash norm\{x\}} $\to \norm{x}$ \quad
          \texttt{\textbackslash BigO\{n\^{}2\}} $\to \BigO{n^2}$
    \item \texttt{\textbackslash inner\{u\}\{v\}} $\to \inner{u}{v}$ \quad
          \texttt{\textbackslash abs\{x\}} $\to \abs{x}$ \quad
          \texttt{\textbackslash mat\{A\}\textbackslash T} $\to \mat{A}\T$
  \end{itemize}
\end{frame}

% =============================================================
\section{Code \& Listings}
% =============================================================

\begin{frame}[fragile]{Code Listings}
  \begin{lstlisting}[language=C, caption={MPI Hello World}]
#include <mpi.h>
#include <stdio.h>

int main(int argc, char** argv) {
    MPI_Init(&argc, &argv);

    int world_size, world_rank;
    MPI_Comm_size(MPI_COMM_WORLD, &world_size);
    MPI_Comm_rank(MPI_COMM_WORLD, &world_rank);

    printf("Hello from rank %d of %d\n", world_rank, world_size);

    MPI_Finalize();
    return 0;
}
  \end{lstlisting}
\end{frame}

% =============================================================
\section{Customisation Guide}
% =============================================================

\begin{frame}{How to Customise}
  \columnsbegin[t]
    \colbegin{0.48\textwidth}
      \blockbegin{Change Colors}
        Edit \texttt{beamercolorthemehorixon.sty}.\\
        Redefine any \texttt{horixon*} color via \texttt{\textbackslash definecolor}.
      \blockend
      \blockbegin{Change Fonts}
        Edit \texttt{beamerfontthemehorixon.sty}.\\
        Swap \texttt{helvet} for any sans-serif package.
      \blockend
    \colend
    \colbegin{0.48\textwidth}
      \blockbegin{Toggle Watermark}
        Pass \texttt{watermark=off} to \texttt{\textbackslash usetheme}\\
        to disable the circuit board overlay.
      \blockend
      \blockbegin{Add Your Logo}
        Pass \texttt{logo=path/to/logo.png}.\\
        It appears in the footer right side.
      \blockend
    \colend
  \columnsend
\end{frame}

% ---- Closing slide ------------------------------------------
\begin{frame}[plain]
  \begin{tikzpicture}[remember picture, overlay]
    \fill[horixonnavy] (current page.south west) rectangle (current page.north east);
    \fill[horixonaccent] (current page.north west) rectangle +(\paperwidth, -0.18cm);
    \fill[horixonaccent] (current page.south west) rectangle +(\paperwidth,  0.12cm);
  \end{tikzpicture}
  \vfill
  \begin{center}
    {\large\color{horixoncream}\bfseries Thank You}\\[0.5em]
    \textcolor{horixonaccent}{\rule{5cm}{0.4pt}}\\[0.4em]
    {\small\color{horixonaccent}\insertauthor}\\
    {\footnotesize\color{horixonaccent}\insertinstitute}\\[0.3em]
    {\footnotesize\color{horixoncream}\insertdate}
  \end{center}
  \vfill
\end{frame}

\end{document}
